% twoopt_swap.tex
% Vizualizácia operácie 2-opt: porovnanie stavu pred a po výmene hrán
% Použitie: \resizebox{0.70\textwidth}{!}{% twoopt_swap.tex
% Vizualizácia operácie 2-opt: porovnanie stavu pred a po výmene hrán
% Použitie: \resizebox{0.70\textwidth}{!}{% twoopt_swap.tex
% Vizualizácia operácie 2-opt: porovnanie stavu pred a po výmene hrán
% Použitie: \resizebox{0.70\textwidth}{!}{% twoopt_swap.tex
% Vizualizácia operácie 2-opt: porovnanie stavu pred a po výmene hrán
% Použitie: \resizebox{0.70\textwidth}{!}{\input{figures/tikz/twoopt_swap}}
\begin{tikzpicture}[
    font=\small,
    vtx/.style={
        circle, draw=black!70, fill=white,
        minimum size=0.55cm, inner sep=1pt, font=\small\bfseries
    },
    %% Hrany trasy
    touredge/.style={draw=blue!65, thick},
    %% Hrany, ktoré sa odstraňujú
    removedge/.style={draw=red!65, thick, dashed},
    %% Hrany, ktoré sa pridávajú
    newedge/.style={draw=green!60!black, very thick},
    %% Obratený úsek
    revedge/.style={draw=blue!40, thick},
    arrow/.style={-{Latex[length=3mm]}, thick, draw=gray!70}
]

%% ══════════════════════════════════════════════════════════════════
%% ĽAVÝ PANEL — pred výmenou
%% ══════════════════════════════════════════════════════════════════
\begin{scope}[xshift=0cm]
    %% Vrcholy (usporiadané lineárne s oblúkmi naznačujúcimi trasu)
    \coordinate (a) at (0.0, 1.5);   % vrchol a
    \coordinate (b) at (1.8, 2.8);   % vrchol b  — hrana (a,b) sa odstráni
    \coordinate (r1) at (2.8, 2.2);  % vnútorný bod úseku
    \coordinate (r2) at (3.2, 1.2);  % vnútorný bod úseku
    \coordinate (c) at (4.0, 2.8);   % vrchol c  — hrana (c,d) sa odstráni
    \coordinate (d) at (5.8, 1.5);   % vrchol d

    %% Hrany trasy mimo (a,b) a (c,d)
    \draw[touredge] (a) -- (b);       % táto sa odstraňuje
    \draw[touredge] (b) -- (r1) -- (r2) -- (c);   % obracajúci sa úsek
    \draw[touredge] (c) -- (d);       % táto sa odstraňuje
    %% Naznačenie pokračovania trasy na krajoch
    \draw[touredge, loosely dotted] (a) -- ++(-0.5, -0.4);
    \draw[touredge, loosely dotted] (d) -- ++(0.5, -0.4);

    %% Červené prerušované hrany (budú sa odstraňovať)
    \draw[removedge] (a) -- (b)
        node[midway, above, font=\scriptsize, red!70] {\(d_{ab}\)};
    \draw[removedge] (c) -- (d)
        node[midway, above, font=\scriptsize, red!70] {\(d_{cd}\)};

    %% Vrcholy
    \node[vtx] at (a) {$a$};
    \node[vtx] at (b) {$b$};
    \node[vtx] at (c) {$c$};
    \node[vtx] at (d) {$d$};

    \node[font=\bfseries\small] at (2.9, -0.15) {Pred výmenou};
    \node[font=\scriptsize, red!70, align=center] at (2.9, -0.70)
        {Odstraňujú sa: \(\{a,b\}\) a \(\{c,d\}\)};
\end{scope}

%% ══════════════════════════════════════════════════════════════════
%% ŠÍPKA MEDZI PANELMI
%% ══════════════════════════════════════════════════════════════════
\draw[arrow] (6.4, 1.5) -- (7.4, 1.5)
    node[midway, above, font=\scriptsize, gray]
        {\(d_{ab}+d_{cd} > d_{ac}+d_{bd}\)};

%% ══════════════════════════════════════════════════════════════════
%% PRAVÝ PANEL — po výmene
%% ══════════════════════════════════════════════════════════════════
\begin{scope}[xshift=7.9cm]
    \coordinate (a) at (0.0, 1.5);
    \coordinate (b) at (1.8, 2.8);
    \coordinate (r1) at (2.8, 2.2);
    \coordinate (r2) at (3.2, 1.2);
    \coordinate (c) at (4.0, 2.8);
    \coordinate (d) at (5.8, 1.5);

    %% Obratený úsek: b→r2→r1→c sa obráti na c→r1→r2→b
    \draw[revedge] (b) -- (r1) -- (r2) -- (c);

    %% Nové hrany (zelené)
    \draw[newedge] (a) -- (c)
        node[midway, above, font=\scriptsize, green!60!black] {\(d_{ac}\)};
    \draw[newedge] (b) -- (d)
        node[midway, above, font=\scriptsize, green!60!black] {\(d_{bd}\)};

    %% Pokračovanie trasy
    \draw[touredge, loosely dotted] (a) -- ++(-0.5, -0.4);
    \draw[touredge, loosely dotted] (d) -- ++(0.5, -0.4);

    %% Vrcholy
    \node[vtx] at (a) {$a$};
    \node[vtx] at (b) {$b$};
    \node[vtx] at (c) {$c$};
    \node[vtx] at (d) {$d$};

    \node[font=\bfseries\small] at (2.9, -0.15) {Po výmene};
    \node[font=\scriptsize, green!60!black, align=center] at (2.9, -0.70)
        {Pridávajú sa: \(\{a,c\}\) a \(\{b,d\}\)};
\end{scope}

\end{tikzpicture}
}
\begin{tikzpicture}[
    font=\small,
    vtx/.style={
        circle, draw=black!70, fill=white,
        minimum size=0.55cm, inner sep=1pt, font=\small\bfseries
    },
    %% Hrany trasy
    touredge/.style={draw=blue!65, thick},
    %% Hrany, ktoré sa odstraňujú
    removedge/.style={draw=red!65, thick, dashed},
    %% Hrany, ktoré sa pridávajú
    newedge/.style={draw=green!60!black, very thick},
    %% Obratený úsek
    revedge/.style={draw=blue!40, thick},
    arrow/.style={-{Latex[length=3mm]}, thick, draw=gray!70}
]

%% ══════════════════════════════════════════════════════════════════
%% ĽAVÝ PANEL — pred výmenou
%% ══════════════════════════════════════════════════════════════════
\begin{scope}[xshift=0cm]
    %% Vrcholy (usporiadané lineárne s oblúkmi naznačujúcimi trasu)
    \coordinate (a) at (0.0, 1.5);   % vrchol a
    \coordinate (b) at (1.8, 2.8);   % vrchol b  — hrana (a,b) sa odstráni
    \coordinate (r1) at (2.8, 2.2);  % vnútorný bod úseku
    \coordinate (r2) at (3.2, 1.2);  % vnútorný bod úseku
    \coordinate (c) at (4.0, 2.8);   % vrchol c  — hrana (c,d) sa odstráni
    \coordinate (d) at (5.8, 1.5);   % vrchol d

    %% Hrany trasy mimo (a,b) a (c,d)
    \draw[touredge] (a) -- (b);       % táto sa odstraňuje
    \draw[touredge] (b) -- (r1) -- (r2) -- (c);   % obracajúci sa úsek
    \draw[touredge] (c) -- (d);       % táto sa odstraňuje
    %% Naznačenie pokračovania trasy na krajoch
    \draw[touredge, loosely dotted] (a) -- ++(-0.5, -0.4);
    \draw[touredge, loosely dotted] (d) -- ++(0.5, -0.4);

    %% Červené prerušované hrany (budú sa odstraňovať)
    \draw[removedge] (a) -- (b)
        node[midway, above, font=\scriptsize, red!70] {\(d_{ab}\)};
    \draw[removedge] (c) -- (d)
        node[midway, above, font=\scriptsize, red!70] {\(d_{cd}\)};

    %% Vrcholy
    \node[vtx] at (a) {$a$};
    \node[vtx] at (b) {$b$};
    \node[vtx] at (c) {$c$};
    \node[vtx] at (d) {$d$};

    \node[font=\bfseries\small] at (2.9, -0.15) {Pred výmenou};
    \node[font=\scriptsize, red!70, align=center] at (2.9, -0.70)
        {Odstraňujú sa: \(\{a,b\}\) a \(\{c,d\}\)};
\end{scope}

%% ══════════════════════════════════════════════════════════════════
%% ŠÍPKA MEDZI PANELMI
%% ══════════════════════════════════════════════════════════════════
\draw[arrow] (6.4, 1.5) -- (7.4, 1.5)
    node[midway, above, font=\scriptsize, gray]
        {\(d_{ab}+d_{cd} > d_{ac}+d_{bd}\)};

%% ══════════════════════════════════════════════════════════════════
%% PRAVÝ PANEL — po výmene
%% ══════════════════════════════════════════════════════════════════
\begin{scope}[xshift=7.9cm]
    \coordinate (a) at (0.0, 1.5);
    \coordinate (b) at (1.8, 2.8);
    \coordinate (r1) at (2.8, 2.2);
    \coordinate (r2) at (3.2, 1.2);
    \coordinate (c) at (4.0, 2.8);
    \coordinate (d) at (5.8, 1.5);

    %% Obratený úsek: b→r2→r1→c sa obráti na c→r1→r2→b
    \draw[revedge] (b) -- (r1) -- (r2) -- (c);

    %% Nové hrany (zelené)
    \draw[newedge] (a) -- (c)
        node[midway, above, font=\scriptsize, green!60!black] {\(d_{ac}\)};
    \draw[newedge] (b) -- (d)
        node[midway, above, font=\scriptsize, green!60!black] {\(d_{bd}\)};

    %% Pokračovanie trasy
    \draw[touredge, loosely dotted] (a) -- ++(-0.5, -0.4);
    \draw[touredge, loosely dotted] (d) -- ++(0.5, -0.4);

    %% Vrcholy
    \node[vtx] at (a) {$a$};
    \node[vtx] at (b) {$b$};
    \node[vtx] at (c) {$c$};
    \node[vtx] at (d) {$d$};

    \node[font=\bfseries\small] at (2.9, -0.15) {Po výmene};
    \node[font=\scriptsize, green!60!black, align=center] at (2.9, -0.70)
        {Pridávajú sa: \(\{a,c\}\) a \(\{b,d\}\)};
\end{scope}

\end{tikzpicture}
}
\begin{tikzpicture}[
    font=\small,
    vtx/.style={
        circle, draw=black!70, fill=white,
        minimum size=0.55cm, inner sep=1pt, font=\small\bfseries
    },
    %% Hrany trasy
    touredge/.style={draw=blue!65, thick},
    %% Hrany, ktoré sa odstraňujú
    removedge/.style={draw=red!65, thick, dashed},
    %% Hrany, ktoré sa pridávajú
    newedge/.style={draw=green!60!black, very thick},
    %% Obratený úsek
    revedge/.style={draw=blue!40, thick},
    arrow/.style={-{Latex[length=3mm]}, thick, draw=gray!70}
]

%% ══════════════════════════════════════════════════════════════════
%% ĽAVÝ PANEL — pred výmenou
%% ══════════════════════════════════════════════════════════════════
\begin{scope}[xshift=0cm]
    %% Vrcholy (usporiadané lineárne s oblúkmi naznačujúcimi trasu)
    \coordinate (a) at (0.0, 1.5);   % vrchol a
    \coordinate (b) at (1.8, 2.8);   % vrchol b  — hrana (a,b) sa odstráni
    \coordinate (r1) at (2.8, 2.2);  % vnútorný bod úseku
    \coordinate (r2) at (3.2, 1.2);  % vnútorný bod úseku
    \coordinate (c) at (4.0, 2.8);   % vrchol c  — hrana (c,d) sa odstráni
    \coordinate (d) at (5.8, 1.5);   % vrchol d

    %% Hrany trasy mimo (a,b) a (c,d)
    \draw[touredge] (a) -- (b);       % táto sa odstraňuje
    \draw[touredge] (b) -- (r1) -- (r2) -- (c);   % obracajúci sa úsek
    \draw[touredge] (c) -- (d);       % táto sa odstraňuje
    %% Naznačenie pokračovania trasy na krajoch
    \draw[touredge, loosely dotted] (a) -- ++(-0.5, -0.4);
    \draw[touredge, loosely dotted] (d) -- ++(0.5, -0.4);

    %% Červené prerušované hrany (budú sa odstraňovať)
    \draw[removedge] (a) -- (b)
        node[midway, above, font=\scriptsize, red!70] {\(d_{ab}\)};
    \draw[removedge] (c) -- (d)
        node[midway, above, font=\scriptsize, red!70] {\(d_{cd}\)};

    %% Vrcholy
    \node[vtx] at (a) {$a$};
    \node[vtx] at (b) {$b$};
    \node[vtx] at (c) {$c$};
    \node[vtx] at (d) {$d$};

    \node[font=\bfseries\small] at (2.9, -0.15) {Pred výmenou};
    \node[font=\scriptsize, red!70, align=center] at (2.9, -0.70)
        {Odstraňujú sa: \(\{a,b\}\) a \(\{c,d\}\)};
\end{scope}

%% ══════════════════════════════════════════════════════════════════
%% ŠÍPKA MEDZI PANELMI
%% ══════════════════════════════════════════════════════════════════
\draw[arrow] (6.4, 1.5) -- (7.4, 1.5)
    node[midway, above, font=\scriptsize, gray]
        {\(d_{ab}+d_{cd} > d_{ac}+d_{bd}\)};

%% ══════════════════════════════════════════════════════════════════
%% PRAVÝ PANEL — po výmene
%% ══════════════════════════════════════════════════════════════════
\begin{scope}[xshift=7.9cm]
    \coordinate (a) at (0.0, 1.5);
    \coordinate (b) at (1.8, 2.8);
    \coordinate (r1) at (2.8, 2.2);
    \coordinate (r2) at (3.2, 1.2);
    \coordinate (c) at (4.0, 2.8);
    \coordinate (d) at (5.8, 1.5);

    %% Obratený úsek: b→r2→r1→c sa obráti na c→r1→r2→b
    \draw[revedge] (b) -- (r1) -- (r2) -- (c);

    %% Nové hrany (zelené)
    \draw[newedge] (a) -- (c)
        node[midway, above, font=\scriptsize, green!60!black] {\(d_{ac}\)};
    \draw[newedge] (b) -- (d)
        node[midway, above, font=\scriptsize, green!60!black] {\(d_{bd}\)};

    %% Pokračovanie trasy
    \draw[touredge, loosely dotted] (a) -- ++(-0.5, -0.4);
    \draw[touredge, loosely dotted] (d) -- ++(0.5, -0.4);

    %% Vrcholy
    \node[vtx] at (a) {$a$};
    \node[vtx] at (b) {$b$};
    \node[vtx] at (c) {$c$};
    \node[vtx] at (d) {$d$};

    \node[font=\bfseries\small] at (2.9, -0.15) {Po výmene};
    \node[font=\scriptsize, green!60!black, align=center] at (2.9, -0.70)
        {Pridávajú sa: \(\{a,c\}\) a \(\{b,d\}\)};
\end{scope}

\end{tikzpicture}
}
\begin{tikzpicture}[
    font=\small,
    vtx/.style={
        circle, draw=black!70, fill=white,
        minimum size=0.55cm, inner sep=1pt, font=\small\bfseries
    },
    %% Hrany trasy
    touredge/.style={draw=blue!65, thick},
    %% Hrany, ktoré sa odstraňujú
    removedge/.style={draw=red!65, thick, dashed},
    %% Hrany, ktoré sa pridávajú
    newedge/.style={draw=green!60!black, very thick},
    %% Obratený úsek
    revedge/.style={draw=blue!40, thick},
    arrow/.style={-{Latex[length=3mm]}, thick, draw=gray!70}
]

%% ══════════════════════════════════════════════════════════════════
%% ĽAVÝ PANEL — pred výmenou
%% ══════════════════════════════════════════════════════════════════
\begin{scope}[xshift=0cm]
    %% Vrcholy (usporiadané lineárne s oblúkmi naznačujúcimi trasu)
    \coordinate (a) at (0.0, 1.5);   % vrchol a
    \coordinate (b) at (1.8, 2.8);   % vrchol b  — hrana (a,b) sa odstráni
    \coordinate (r1) at (2.8, 2.2);  % vnútorný bod úseku
    \coordinate (r2) at (3.2, 1.2);  % vnútorný bod úseku
    \coordinate (c) at (4.0, 2.8);   % vrchol c  — hrana (c,d) sa odstráni
    \coordinate (d) at (5.8, 1.5);   % vrchol d

    %% Hrany trasy mimo (a,b) a (c,d)
    \draw[touredge] (a) -- (b);       % táto sa odstraňuje
    \draw[touredge] (b) -- (r1) -- (r2) -- (c);   % obracajúci sa úsek
    \draw[touredge] (c) -- (d);       % táto sa odstraňuje
    %% Naznačenie pokračovania trasy na krajoch
    \draw[touredge, loosely dotted] (a) -- ++(-0.5, -0.4);
    \draw[touredge, loosely dotted] (d) -- ++(0.5, -0.4);

    %% Červené prerušované hrany (budú sa odstraňovať)
    \draw[removedge] (a) -- (b)
        node[midway, above, font=\scriptsize, red!70] {\(d_{ab}\)};
    \draw[removedge] (c) -- (d)
        node[midway, above, font=\scriptsize, red!70] {\(d_{cd}\)};

    %% Vrcholy
    \node[vtx] at (a) {$a$};
    \node[vtx] at (b) {$b$};
    \node[vtx] at (c) {$c$};
    \node[vtx] at (d) {$d$};

    \node[font=\bfseries\small] at (2.9, -0.15) {Pred výmenou};
    \node[font=\scriptsize, red!70, align=center] at (2.9, -0.70)
        {Odstraňujú sa: \(\{a,b\}\) a \(\{c,d\}\)};
\end{scope}

%% ══════════════════════════════════════════════════════════════════
%% ŠÍPKA MEDZI PANELMI
%% ══════════════════════════════════════════════════════════════════
\draw[arrow] (6.4, 1.5) -- (7.4, 1.5)
    node[midway, above, font=\scriptsize, gray]
        {\(d_{ab}+d_{cd} > d_{ac}+d_{bd}\)};

%% ══════════════════════════════════════════════════════════════════
%% PRAVÝ PANEL — po výmene
%% ══════════════════════════════════════════════════════════════════
\begin{scope}[xshift=7.9cm]
    \coordinate (a) at (0.0, 1.5);
    \coordinate (b) at (1.8, 2.8);
    \coordinate (r1) at (2.8, 2.2);
    \coordinate (r2) at (3.2, 1.2);
    \coordinate (c) at (4.0, 2.8);
    \coordinate (d) at (5.8, 1.5);

    %% Obratený úsek: b→r2→r1→c sa obráti na c→r1→r2→b
    \draw[revedge] (b) -- (r1) -- (r2) -- (c);

    %% Nové hrany (zelené)
    \draw[newedge] (a) -- (c)
        node[midway, above, font=\scriptsize, green!60!black] {\(d_{ac}\)};
    \draw[newedge] (b) -- (d)
        node[midway, above, font=\scriptsize, green!60!black] {\(d_{bd}\)};

    %% Pokračovanie trasy
    \draw[touredge, loosely dotted] (a) -- ++(-0.5, -0.4);
    \draw[touredge, loosely dotted] (d) -- ++(0.5, -0.4);

    %% Vrcholy
    \node[vtx] at (a) {$a$};
    \node[vtx] at (b) {$b$};
    \node[vtx] at (c) {$c$};
    \node[vtx] at (d) {$d$};

    \node[font=\bfseries\small] at (2.9, -0.15) {Po výmene};
    \node[font=\scriptsize, green!60!black, align=center] at (2.9, -0.70)
        {Pridávajú sa: \(\{a,c\}\) a \(\{b,d\}\)};
\end{scope}

\end{tikzpicture}
